\section*{Chapter \ref{ch:ml:probabilistic-models-and-uncertainty} --- Probabilistic Models and Uncertainty}

\begin{solution}{exo:ml:probabilistic-models-and-uncertainty:1}
$y - q = 1\%$: loss $0.95\times1\% = 0.95\%$. With $y = 1\%$, $y - q = -1\%$: loss $(0.95 - 1)(-1\%) = 0.05\%$. The
asymmetry is what makes the minimiser the 95\% quantile.
\end{solution}

\begin{solution}{exo:ml:probabilistic-models-and-uncertainty:2}
$\lceil251\times0.9\rceil = 226$th smallest of 250; $\lceil5\,041\times0.9\rceil = 4\,537$th of 5\,040.
\end{solution}

\begin{solution}{exo:ml:probabilistic-models-and-uncertainty:3}
Weights $\propto\mu/\sigma^2$: $5/80^2 : 5/250^2 = 250^2/80^2 = 9.8 : 1$.
\end{solution}

\begin{solution}{exo:ml:probabilistic-models-and-uncertainty:4}
The raw interval has a fixed width set on the calm calibration year; when every shock is 1.6 times larger the same width
covers far fewer returns (74.0\%). The normalised score divides by the predicted scale, which rises with the EWMA
volatility after the break, so the width follows the new regime, with a lag (86.9\%).
\end{solution}

\begin{solution}{exo:ml:probabilistic-models-and-uncertainty:5}
Almost all the predictive variance is the return's own randomness; more data would narrow the forecast very little. The
epistemic part is larger where the model extrapolates: inputs outside the training range, new instruments, small
samples.
\end{solution}

\begin{solution}{exo:ml:probabilistic-models-and-uncertainty:6}
The standard error of an annualised Sharpe ratio over three years is about $\sqrt{1/3} = 0.58$; a difference of 0.3 is
noise. On twenty-year samples and four seeds the same rule raises the Sharpe ratio (1.72 to 1.91). Test sizing rules on
long samples or many seeds.
\end{solution}

\begin{solution}{exo:ml:probabilistic-models-and-uncertainty:7}
\texttt{conformal(gamma=0.02)}: coverage after the break 90.0\% (90.1\% with $\gamma = 0.005$). A larger $\gamma$ reacts faster
to a run of misses but moves the threshold more from day to day, widening and narrowing the band on noise; a smaller one
is smoother and slower after a break.
\end{solution}

\begin{solution}{exo:ml:probabilistic-models-and-uncertainty:8}
$\ell_\alpha(y, q) = (\mathbf 1\{y < q\} - \alpha)(q - y)$. Substituting $q = F^{-1}(\alpha)$, $\alpha = F(z)$, $dz =
dF^{-1}(\alpha)$ and integrating by parts, $2\int_0^1\ell_\alpha(y, F^{-1}(\alpha))\,d\alpha = 2\int(\mathbf 1\{y < z\} -
F(z))(z - y)\,dF(z) = \int(F(z) - \mathbf 1\{z\ge y\})^2dz$ (Gneiting and Raftery, 2007, give the identity).
\end{solution}

\begin{solution}{pb:ml:probabilistic-models-and-uncertainty:1}
\begin{enumerate}
\item A proper score's expectation is minimised by the true distribution, so no forecaster gains by shading; the
log-likelihood, the \omterm{def:ml:probabilistic-models-and-uncertainty:crps}{pinball loss} and the CRPS are proper.
\item CRPS 78.8--79.1 basis points for all, coverage 89.4--91.1\%.
\item CRPS 135.3--136.3, coverage 85.0--87.6\% (the truth 91.3\%).
\item The score is dominated by the volatility, which all of them forecast from the same persistent features; the mean
forecast's signal is a few hundredths of the variance.
\item 0.09\%.
\item The split into aleatoric and epistemic variance, and a slightly better calibrated interval (91.1\% and 87.6\% against
90.6\% and 87.0\% for one seed).
\item Skewness and fat tails through two components (here with coverage 89.7\% and 86.1\%).
\item Out of the training range, for new instruments, and with few labels.
\item Raw: 90.7\% and 74.0\%; normalised: 90.3\% and 86.9\%.
\item 89.9\% and 90.1\%; it gives up the finite-sample guarantee for a long-run error rate under drift.
\item Exchangeability of calibration and test scores.
\item From the speed of the drifts expected and the tolerance for a wobbling width: small (0.005) for slow changes.
\item Over three years Sharpe ratios have a standard error near 0.6, larger than the effect.
\item Reading 1.72; over the EWMA variance 1.90, over the network's 1.91, over the true 1.93. True mean 2.94, over the true
variance 3.31.
\item Scaling reweights the signal's expected profits; the more of the mean is signal rather than noise, the more there is
to reweight.
\item \emph{Named result}: dividing the drift reading by the network's predicted variance raises the Sharpe ratio from 1.72
to 1.91 (1.93 with the true variance); after the regime change split conformal intervals cover 74.0\% (raw) or 86.9\%
(normalised) of the nominal 90\%, adaptive conformal 90.1\%.
\item The predictive standard deviation or an interval, and its recent coverage.
\item Track the interval's hit rate and the normalised errors' distribution with the drift detectors of chapter~12 and the
alarms of chapter~27.
\item Too-narrow tails: a Gaussian model under-covers extremes; use quantiles, a mixture, or conformal intervals.
\item A distribution, of which the point forecast is one summary.
\end{enumerate}
\end{solution}

\begin{solution}{iq:ml:probabilistic-models-and-uncertainty:1}
A score whose expected value is minimised by forecasting the true distribution. Training with it makes honest
distributions the optimum; evaluating with it lets different forecasts be compared without rewarding over- or
under-confidence.
\iqlookfor{the definition and the incentive argument.}
\end{solution}

\begin{solution}{iq:ml:probabilistic-models-and-uncertainty:2}
Fit quantile objectives (\omterm{def:ml:probabilistic-models-and-uncertainty:crps}{pinball loss}) at the interval's levels, or a mean model plus a model of the absolute residual,
and wrap either in conformal calibration on held-out data (conformalised \omterm{def:ml:probabilistic-models-and-uncertainty:quantile}{quantile regression}).
\iqlookfor{quantile losses and calibration on held-out data.}
\end{solution}

\begin{solution}{iq:ml:probabilistic-models-and-uncertainty:3}
Compute scores on a calibration set; the threshold is their $\lceil(n+1)(1-\alpha)\rceil$-th smallest; intervals are the
values whose score is below it; coverage is at least $1 - \alpha$ if calibration and new scores are exchangeable. It
fails when markets change regime (volatility, liquidity), since the calibration scores no longer describe the new ones;
use normalised scores and adaptive updates.
\iqlookfor{the construction, the exchangeability condition, and the remedy.}
\end{solution}

\begin{solution}{iq:ml:probabilistic-models-and-uncertainty:4}
In proportion to the mean over the variance (Kelly or mean-variance), possibly shrunk; test on long histories or many
simulated paths, since the Sharpe ratio differences between rules are small relative to their standard errors.
\iqlookfor{mean over variance, and a test with enough power.}
\end{solution}

\begin{solution}{iq:ml:probabilistic-models-and-uncertainty:5}
Aleatoric: the target's noise given the inputs, estimated by a predicted variance. Epistemic: uncertainty about the
model, estimated by the disagreement of an ensemble (or a posterior). The first stays with more data, the second shrinks.
\iqlookfor{the definitions and an estimator for each.}
\end{solution}

\begin{solution}{iq:ml:probabilistic-models-and-uncertainty:6}
$\frac{d}{dq}\E\ell_\alpha(Y, q) = G(q) - \alpha$ for continuous $G$: negative below the $\alpha$-quantile, positive above, so
the minimum is at $G^{-1}(\alpha)$.
\iqlookfor{the derivative and the sign argument.}
\end{solution}
